\documentclass[12pt]{article}

\usepackage[utf8]{inputenc}
\usepackage{amsmath}
\usepackage{enumitem}
\usepackage{xcolor}
\usepackage{geometry}

\geometry{margin=1in}

\title{Prompting Techniques}
\author{Datta Punjare}
\date{}

\begin{document}

\maketitle

\section{One-Shot Prompting}

\begin{itemize}
    \item One-shot prompting means giving the model one example of the task, so it can understand the pattern and produce a similar answer.
\end{itemize}

\subsection*{Input}
\begin{verbatim}
Hello -> Bonjour
Now translate: GoodBye
\end{verbatim}

\subsection*{Output}
\begin{verbatim}
Goodbye -> Au revoir
\end{verbatim}

\subsection*{Observation}
Only one example is given as a demonstration at the start of the chat. The model observes the given example, understands the pattern, and applies the same pattern to the new input.

% ------------------------------------------------

\section{Few-Shot Prompting}

\begin{itemize}
    \item Few-shot prompting means giving the model multiple examples of a task so that it can understand the pattern and apply it to a new input.
\end{itemize}

\subsection*{Input}
\begin{verbatim}
Hello -> Bonjour
Thanks -> Merci
Now translate: GoodBye
\end{verbatim}

\subsection*{Output}
\begin{verbatim}
GoodBye -> Au revoir
\end{verbatim}

\subsection*{Observation}
Provide multiple examples to the model. By observing these examples, the model understands the pattern better and applies it to the new input.

% ------------------------------------------------

\section{Zero-Shot Prompting}

\begin{itemize}
    \item Zero-shot prompting means giving the model a task without providing any examples. It gives only an instruction or question and tells the model the steps to achieve the answer.
\end{itemize}

\subsection*{Input}
\begin{verbatim}
If a train travels 60 km in 1.5 hours, what is its speed?
\end{verbatim}

\subsection*{Output}
\begin{verbatim}
Speed = Distance / Time
60 / 1.5 = 40
Answer: 40 km/h
\end{verbatim}

\subsection*{Observation}
Zero-shot prompting does not provide any example to the AI. It directly asks a question to the model. The model answers the question using its pre-trained knowledge.

% ------------------------------------------------

\section{Natural Language Query (Direct Answer)}

\subsection*{Input}
\begin{verbatim}
What's the weather in Nashik?
\end{verbatim}

\subsection*{Output}
\begin{verbatim}
Right now in Nashik, it is around 28 degrees C with rain.
The weather is breezy, with more showers possible this afternoon.
\end{verbatim}

\subsection*{Observation}
The model did not generate the weather details itself. Instead, it took the information from a live weather API connected in the background.

% ------------------------------------------------

\section{Code Generation Request (Function Simulation)}

\subsection*{Input}
\begin{verbatim}
Write a Python function get_weather(city)
and show how it would call it for Nashik.
\end{verbatim}

\subsection*{Output}
\begin{verbatim}
Weather information for Nashik
\end{verbatim}

% ------------------------------------------------

\section{Call the Function for Nashik}

\subsection*{Input}
\begin{verbatim}
Write a Python function get_weather(nashik)
and show how it would call it for Nashik.
\end{verbatim}

\subsection*{Output}
\begin{verbatim}
Getting weather for Nashik
\end{verbatim}

\end{document}
